% !TEX program = pdflatex
\documentclass[11pt,a4paper]{article}

% =========================================================================
% INFORME DE VALIDACIÓN DEL PRODUCTO
% Proyecto: CasaScore (AlquilaPuno)
% Gestión de Proyectos TICS (EST315) -- Semestre 2026-II -- Team Sobrevivientes
% =========================================================================

% --- PAQUETES DE CONFIGURACIÓN ---
\usepackage[utf8]{inputenc}
\usepackage[T1]{fontenc}
\usepackage[spanish,es-tabla]{babel}
% Compatibilidad con babel-spanish reciente (flechas de TikZ y signo de porcentaje)
\AtBeginDocument{\shorthandoff{<>}\spanishplainpercent}
\usepackage{geometry}
\geometry{
    top=3.6cm,
    bottom=2.3cm,
    left=2.2cm,
    right=2.2cm,
    % PARTE MODIFICABLE: headheight debe ser mayor que la altura del logo más alto (2.56cm = 73pt).
    headheight=80pt,
    headsep=15pt,
    footskip=20pt
}

\usepackage{graphicx}
\usepackage{fancyhdr}
\usepackage[table]{xcolor}
\usepackage{titlesec}
\usepackage{booktabs}
\usepackage{tabularx}
\usepackage{array}
\usepackage{amsmath}
\usepackage{enumitem}
\usepackage{microtype}
\usepackage{amssymb}
% --- RECURSOS VISUALES: GRÁFICOS Y CÓDIGO DE COLORES ---
\usepackage{tikz}
\usetikzlibrary{positioning,calc}
\definecolor{GrisFila}{RGB}{246, 248, 251}
\definecolor{CuadCerca}{RGB}{150, 30, 30}
\definecolor{CuadSatisf}{RGB}{180, 130, 20}
\definecolor{CuadInfo}{RGB}{10, 45, 90}
\definecolor{CuadMonit}{RGB}{110, 115, 125}
\newcommand{\etiqueta}[2]{{\setlength{\fboxsep}{2pt}\colorbox{#1}{\textcolor{white}{\scriptsize\bfseries #2}}}}
\newcommand{\cerca}{\etiqueta{CuadCerca}{Gestionar de cerca}}
\newcommand{\satisf}{\etiqueta{CuadSatisf}{Mantener satisfecho}}
\newcommand{\informado}{\etiqueta{CuadInfo}{Mantener informado}}
\newcommand{\monit}{\etiqueta{CuadMonit}{Monitorear}}
\usepackage{needspace}
\usepackage{longtable}
\definecolor{Rama1}{RGB}{10, 45, 90}
\definecolor{Rama2}{RGB}{40, 95, 150}
\definecolor{Rama3}{RGB}{180, 130, 20}
\definecolor{Rama4}{RGB}{30, 110, 60}
\definecolor{Rama5}{RGB}{110, 60, 140}
\definecolor{RaciR}{RGB}{150, 30, 30}
\definecolor{RaciA}{RGB}{10, 45, 90}
\definecolor{RaciC}{RGB}{180, 130, 20}
\definecolor{RaciI}{RGB}{140, 145, 155}
\definecolor{NivAlto}{RGB}{30, 110, 60}
\definecolor{NivMedio}{RGB}{180, 130, 20}
\definecolor{NivBajo}{RGB}{150, 30, 30}
\usetikzlibrary{patterns}
\definecolor{SevBaja}{RGB}{30, 110, 60}
\definecolor{SevMedia}{RGB}{180, 130, 20}
\definecolor{SevAlta}{RGB}{150, 30, 30}
\usetikzlibrary{arrows.meta,shapes.geometric,fit,backgrounds}
\usepackage{listings}
\lstset{basicstyle=\ttfamily\scriptsize, frame=single, rulecolor=\color{AzulUNA!40}, backgroundcolor=\color{GrisFila}, columns=fullflexible, keepspaces=true, breaklines=true, xleftmargin=4pt, xrightmargin=4pt, literate={á}{{\'a}}1 {é}{{\'e}}1 {í}{{\'i}}1 {ó}{{\'o}}1 {ú}{{\'u}}1 {ñ}{{\~n}}1 {Á}{{\'A}}1}

\usepackage{hyperref}

% Logos: funciona a 2 o 3 niveles de profundidad dentro de docs/latex/
\graphicspath{{../../assets/img/}{../../../assets/img/}}

% --- PALETA INSTITUCIONAL FIEI / UNA PUNO ---
\definecolor{AzulUNA}{RGB}{10, 45, 90}
\definecolor{DoradoFIEI}{RGB}{180, 130, 20}
\definecolor{GrisOscuro}{RGB}{50, 50, 50}
\definecolor{AzulClaro}{RGB}{228, 236, 246}

% --- CONFIGURACIÓN DE LISTAS ---
\setlist{topsep=2pt, itemsep=2pt, parsep=1pt}

% --- COLUMNAS DE TABLA ---
\newcolumntype{L}{>{\raggedright\arraybackslash}X}
\newcolumntype{P}[1]{>{\raggedright\arraybackslash}p{#1}}
\newcolumntype{C}[1]{>{\centering\arraybackslash}p{#1}}
\newcolumntype{Q}[1]{>{\raggedleft\arraybackslash}p{#1}}
\newcolumntype{R}{>{\raggedleft\arraybackslash}X}
\renewcommand{\arraystretch}{1.25}

% --- FORMATO DE ENLACES ---
\hypersetup{
    colorlinks=true,
    linkcolor=AzulUNA,
    urlcolor=AzulUNA,
    citecolor=AzulUNA,
    pdftitle={Informe de Validación del Producto -- CasaScore (AlquilaPuno)},
    pdfauthor={Team Sobrevivientes}
}

% --- FORMATO DE TÍTULOS ---
\titleformat{\section}{\color{AzulUNA}\normalfont\large\bfseries}{\thesection.}{0.7em}{}[\color{AzulUNA}\titlerule]
\titleformat{\subsection}{\color{AzulUNA}\normalfont\normalsize\bfseries}{\thesubsection.}{0.5em}{}
\titlespacing*{\section}{0pt}{9pt}{4pt}
\titlespacing*{\subsection}{0pt}{6pt}{2pt}

% --- ENCABEZADO Y PIE DE PÁGINA (fancyhdr) ---
\pagestyle{fancy}
\fancyhf{}

% ===== TAMAÑO DE LOS LOGOS (PARTE MODIFICABLE) =====
% logo_fiei.png  : height=1.45cm (lado izquierdo)
% logo_equipo.png: height=2.56cm (lado derecho; 25 % más grande porque la imagen tiene más borde blanco)
% Para cambiar el tamaño, modifique solo el valor height=... de cada \includegraphics.
% Si el logo más alto supera 2.6cm, aumente también headheight y top en la configuración de página.
\fancyhead[L]{%
    \includegraphics[height=1.45cm,keepaspectratio]{logo_fiei.png}%
}
\fancyhead[C]{%
    \scriptsize \textbf{\textcolor{AzulUNA}{UNIVERSIDAD NACIONAL DEL ALTIPLANO DE PUNO}}\\[2pt]
    \tiny \textcolor{GrisOscuro}{FACULTAD DE INGENIERÍA ESTADÍSTICA E INFORMÁTICA}\\[2pt]
    \scriptsize \textbf{\textcolor{AzulUNA}{Gestión de Proyectos TICS (EST315) \textbar\ Team Sobrevivientes}}\\[2pt]
    \footnotesize \textbf{\textcolor{DoradoFIEI}{INFORME DE VALIDACIÓN DEL PRODUCTO}}%
}
\fancyhead[R]{%
    \includegraphics[height=2.56cm,keepaspectratio]{logo_equipo.png}%
}

\renewcommand{\headrulewidth}{0pt}
\renewcommand{\headrule}{}

\fancyfoot[L]{\scriptsize \textcolor{GrisOscuro}{CasaScore (AlquilaPuno)}}
\fancyfoot[C]{\small \thepage}
\fancyfoot[R]{\scriptsize \textcolor{GrisOscuro}{Puno, Perú -- Semestre 2026-II}}
\renewcommand{\footrulewidth}{0.4pt}

\begin{document}

\begin{center}
    {\Large \textbf{\textcolor{AzulUNA}{INFORME DE VALIDACIÓN DEL PRODUCTO}}}\\[3pt]
    {\normalsize \textbf{CasaScore (AlquilaPuno) -- Prototipo funcional del MVP de la Unidad I}}\\[4pt]
    \footnotesize \textbf{Fecha:} viernes 09/10/2026 (semana 7) \quad \textbar\quad \textbf{Versión:} 1.1 \quad \textbar\quad \textbf{Fase:} Implementación
\end{center}

\vspace{0.1cm}

\section{Datos Generales}
\noindent\begin{tabularx}{\textwidth}{@{}P{4.2cm} L@{}}
\toprule
\textbf{Proyecto:} & CasaScore (aplicación AlquilaPuno) \\
\textbf{Responsable:} & Eliseo Tarqui Ajahuana -- Especialista Técnico / Arquitectura \\
\textbf{Proceso PM4DEV:} & Gestión de la calidad (control del producto) \\
\textbf{Producto evaluado:} & Aplicación Android con el motor CASA SCORE, el ranking explicado y la publicación de cuartos \\
\textbf{Insumos:} & Especificación de requerimientos, plan de calidad y arquitectura técnica \\
\bottomrule
\end{tabularx}

\section{Alcance del Prototipo}
El prototipo valida el núcleo del producto: que el motor ordene los cuartos de forma correcta y explicada. Las funciones que dependen de la nube se construyen para el MVP final.

\vspace{0.1cm}
{\small
\rowcolors{2}{white}{GrisFila}
\noindent\begin{tabularx}{\textwidth}{@{}C{1.2cm} P{3.3cm} C{2.1cm} L@{}}
\toprule
\rowcolor{AzulClaro}
\textbf{ID} & \textbf{Requerimiento} & \textbf{Estado} & \textbf{Detalle} \\
\midrule
RF-01 & Registro e inicio de sesión & \etiqueta{SevAlta}{MVP final} & Se implementa con Supabase en las semanas 10 a 12. \\
RF-02 & Publicación de cuartos & \etiqueta{SevMedia}{Parcial} & Formulario con 4 fotos obligatorias, servicios, condiciones y área. La ubicación en el mapa y el cálculo automático de la distancia quedan para el MVP final. \\
RF-03 & Punto de referencia & \etiqueta{SevAlta}{MVP final} & El prototipo usa la distancia y el tiempo a la UNA Puno ingresados por el propietario. \\
RF-04 & Prioridades del arrendatario & \etiqueta{SevMedia}{Parcial} & 4 pesos con 4 perfiles predefinidos: precio, cercanía, seguridad y equilibrado. Falta el peso del baño. \\
RF-05 & Normalización 0--100 & \etiqueta{Rama4}{Implementado} & Precio, cercanía, seguridad y servicios. \\
RF-06 & CASA SCORE y ranking & \etiqueta{Rama4}{Implementado} & Suma ponderada, bonificaciones, penalizaciones y límite de 0 a 100. \\
RF-07 & Explicación del puntaje & \etiqueta{Rama4}{Implementado} & Puntaje por criterio y lista de bonificaciones y penalizaciones. \\
RF-08 & Perfil y mis cuartos & \etiqueta{SevAlta}{MVP final} & Se implementa con las cuentas de usuario. \\
RF-09 & Datos en la nube & \etiqueta{SevAlta}{MVP final} & El prototipo guarda los cuartos en SQLite dentro del celular. \\
\bottomrule
\end{tabularx}
}

\Needspace{14\baselineskip}
\section{Arquitectura del Prototipo}
{\small
\rowcolors{2}{white}{GrisFila}
\noindent\begin{tabularx}{\textwidth}{@{}P{4.3cm} L@{}}
\toprule
\rowcolor{AzulClaro}
\textbf{Componente} & \textbf{Implementación} \\
\midrule
Plataforma & Expo SDK 57, React Native 0.86, React 19.2 y TypeScript 6.0 \\
Pantalla principal (\texttt{App.tsx}) & Ranking de cuartos, selector de 4 perfiles de prioridades y detalle con el desglose del puntaje \\
Publicación (\texttt{AddPropertyModal.tsx}) & Formulario con datos, servicios, condiciones y 4 fotos con \texttt{expo-image-picker} \\
Motor (\texttt{utils/scoring.ts}) & Normalización, suma ponderada, bonificaciones, penalizaciones y notas explicativas, separado de las pantallas \\
Datos (\texttt{data/database.ts}) & Tabla de cuartos en SQLite (\texttt{expo-sqlite}) con 5 cuartos de prueba de Puno precargados \\
\bottomrule
\end{tabularx}
}

\section{Reglas Implementadas frente a la Especificación}
{\small
\rowcolors{2}{white}{GrisFila}
\noindent\begin{tabularx}{\textwidth}{@{}P{2.4cm} L L@{}}
\toprule
\rowcolor{AzulClaro}
\textbf{Criterio} & \textbf{En el prototipo} & \textbf{Cambio para el MVP final} \\
\midrule
Precio & Normalización entre S/ 140 y S/ 850. & Sin cambio. \\
Cercanía & 50\,\% distancia y 50\,\% tiempo a la UNA Puno, ingresados por el propietario. & Distancia y tiempo calculados al punto de referencia del arrendatario. \\
Seguridad & Calificación de 1 a 5 ingresada por el propietario. & 1 punto base más 1 por cada característica declarada con foto. \\
Servicios 24 h & Agua 40, luz 30 e internet 30. & Sin cambio. \\
Baño & No implementado. & Privado 100 y compartido 0, con su propio peso. \\
Bonificaciones y penalizaciones & Las 3 bonificaciones y las 6 penalizaciones de la especificación. & Revisar el tope de 100 (sección 8). \\
\bottomrule
\end{tabularx}
}

\Needspace{14\baselineskip}
\section{Casos de Prueba del Motor}
Se ejecutó el motor de la aplicación (\texttt{scoring.ts}) con los valores de cada caso del plan de calidad. Los casos manuales en celular (registro, mapa, nube y rendimiento en el equipo) se ejecutan con el MVP final.

\vspace{0.1cm}
{\footnotesize
\setlength{\tabcolsep}{3.5pt}
\rowcolors{2}{white}{GrisFila}
\noindent\begin{tabularx}{\textwidth}{@{}C{1.0cm} P{3.6cm} L L C{1.6cm}@{}}
\toprule
\rowcolor{AzulClaro}
\textbf{Caso} & \textbf{Entrada} & \textbf{Esperado} & \textbf{Obtenido} & \textbf{Resultado} \\
\midrule
CP-01 & Precio de S/ 140, S/ 850 y S/ 495 & 100, 0 y 50 & 100, 0 y 50 & \etiqueta{Rama4}{Aprobado} \\
CP-02 & Cercanía: 50 m y 2 min; 3,000 m y 40 min & 100 y 0 & 100 y 0 & \etiqueta{Rama4}{Aprobado} \\
CP-03 & Seguridad de nivel 1 y nivel 5 & 20 y 100 & 20 y 100 & \etiqueta{Rama4}{Aprobado} \\
CP-04 & Servicios: agua y luz, sin internet & 70 & 70 & \etiqueta{Rama4}{Aprobado} \\
CP-05 & Bonificaciones, penalizaciones y límite & Luz $+8$; llegada 20:00 $-20$; no declarado $-25$; resultado entre 0 y 100 & Igual; el cuarto Bellavista baja a 8 y el de Puerta 1 se limita a 100 & \etiqueta{Rama4}{Aprobado} \\
CP-06 & Solo el peso del precio (100/0/0/0) & Puntaje base igual al del precio (86) & 86 & \etiqueta{Rama4}{Aprobado} \\
CP-07 & Orden de los 5 cuartos de prueba & Puerta 1, Tesista, Laykakota, Huáscar, Bellavista & Mismo orden & \etiqueta{Rama4}{Aprobado} \\
\bottomrule
\end{tabularx}
}

\Needspace{16\baselineskip}
\section{Ranking de los Cuartos de Prueba}
Resultado con el perfil por defecto (precio 40, cercanía 30, seguridad 20 y servicios 10). B: bonificaciones; P: penalizaciones; Score: CASA SCORE final.

\vspace{0.1cm}
{\footnotesize
\setlength{\tabcolsep}{3pt}
\rowcolors{2}{white}{GrisFila}
\noindent\begin{tabularx}{\textwidth}{@{}C{0.6cm} L Q{0.9cm} C{1.9cm} C{0.75cm} C{0.75cm} C{0.75cm} C{0.75cm} C{0.8cm} C{0.65cm} C{0.65cm} C{1.1cm}@{}}
\toprule
\rowcolor{AzulClaro}
\textbf{N.\textsuperscript{o}} & \textbf{Cuarto} & \textbf{S/} & \textbf{Distancia} & \textbf{Prec.} & \textbf{Cerc.} & \textbf{Seg.} & \textbf{Serv.} & \textbf{Base} & \textbf{B} & \textbf{P} & \textbf{Score} \\
\midrule
1 & Hab. Puerta 1 (Av. Floral) & 240 & 120 m / 2 min & 86 & 99 & 90 & 100 & 92 & $+$13 & $-$0 & \etiqueta{Rama4}{\,100\,} \\
2 & Hab. Tesista (Sesquicentenario) & 320 & 380 m / 5 min & 75 & 91 & 96 & 100 & 87 & $+$12 & $-$6 & \etiqueta{Rama4}{\,93\,} \\
3 & Minidepartamento Laykakota & 480 & 600 m / 8 min & 52 & 83 & 84 & 100 & 73 & $+$17 & $-$0 & \etiqueta{Rama4}{\,90\,} \\
4 & Cuarto Huáscar (2.\textsuperscript{o} piso) & 200 & 800 m / 12 min & 92 & 75 & 70 & 100 & 83 & $+$8 & $-$24 & \etiqueta{SevMedia}{\,67\,} \\
5 & Cuarto Bellavista (Cahuide) & 160 & 1,100 m / 16 min & 97 & 64 & 60 & 30 & 73 & $+$0 & $-$65 & \etiqueta{SevAlta}{\,8\,} \\
\bottomrule
\end{tabularx}
}

\vspace{0.15cm}
\begin{itemize}[leftmargin=*]
    \item \textbf{Puerta 1} suma 92 de base y 13 de bonificación (105): el límite lo deja en 100.
    \item \textbf{Tesista} supera a \textbf{Laykakota} por su precio menor y su mejor seguridad, aunque pierde 6 puntos por la hora de llegada de las 23:00.
    \item \textbf{Bellavista} tiene el mejor precio, pero baja a 8 por llegada hasta las 20:00, sin visitas, sin cocina y condiciones no declaradas ($-65$).
\end{itemize}

\Needspace{12\baselineskip}
\section{Sensibilidad ante las Prioridades}
CASA SCORE de cada cuarto según el perfil elegido. En negrita, el segundo lugar en cada perfil.

\vspace{0.1cm}
{\footnotesize
\rowcolors{2}{white}{GrisFila}
\noindent\begin{tabularx}{\textwidth}{@{}L C{1.5cm} C{1.5cm} C{1.6cm} C{1.5cm} C{1.6cm}@{}}
\toprule
\rowcolor{AzulClaro}
\textbf{Perfil (pesos de precio, cercanía, seguridad y servicios)} & \textbf{Puerta 1} & \textbf{Tesista} & \textbf{Laykakota} & \textbf{Huáscar} & \textbf{Bellavista} \\
\midrule
Por defecto (40/30/20/10) & 100 & \textbf{93} & 90 & 67 & 8 \\
Precio (60/20/10/10) & 100 & \textbf{89} & 83 & 71 & 15 \\
Cercanía (20/60/10/10) & 100 & 95 & \textbf{96} & 64 & 2 \\
Seguridad (20/20/50/10) & 100 & \textbf{97} & 96 & 62 & 0 \\
Equilibrado (25/25/25/25) & 100 & \textbf{97} & \textbf{97} & 68 & 0 \\
\bottomrule
\end{tabularx}
}

\vspace{0.15cm}
\noindent El primer y los dos últimos lugares no cambian. El segundo lugar sí depende de las prioridades: con el perfil de cercanía gana Laykakota, y con el de precio gana Tesista. Esto confirma que el ranking responde a lo que cada arrendatario valora.

\Needspace{22\baselineskip}
\section{Hallazgos y Correcciones}
{\small
\rowcolors{2}{white}{GrisFila}
\noindent\begin{tabularx}{\textwidth}{@{}C{0.7cm} L L C{1.5cm}@{}}
\toprule
\rowcolor{AzulClaro}
\textbf{N.\textsuperscript{o}} & \textbf{Hallazgo} & \textbf{Corrección} & \textbf{Semana} \\
\midrule
1 & Las bonificaciones llevan al mejor cuarto por encima de 100 y el límite le quita diferencia con el segundo. & Evaluar con el equipo una escala de bonificaciones menor o aplicarlas solo hasta 100 antes de restar penalizaciones. & 11 \\
2 & La distancia la ingresa el propietario y siempre es a la UNA Puno. & Cálculo automático al punto de referencia del arrendatario. & 11 -- 12 \\
3 & La seguridad es una calificación declarada sin sustento. & Características con foto, según la especificación. & 11 -- 12 \\
4 & El tipo de baño no influye en el puntaje. & Agregar el quinto criterio y su peso. & 11 -- 12 \\
5 & Las pruebas del motor se ejecutan a mano. & Pruebas automáticas con Jest en cada Pull Request. & 8 -- 10 \\
6 & Los datos solo existen en el celular. & Migración a Supabase con cuentas y seguridad por fila. & 10 -- 12 \\
\bottomrule
\end{tabularx}
}

\section{Estado frente al Cronograma}
\begin{itemize}[leftmargin=*]
    \item El prototipo se entregó el 09/10/2026, tres días antes de la fecha de la línea base para el MVP de la Unidad I (12/10/2026).
    \item Al calcular 500 cuartos y ordenarlos en una computadora, el motor tardó menos de 5 ms; la medición en celular se hace con el MVP final.
    \item Las actividades de gestión del proyecto mantienen un atraso de dos semanas, registrado en las actas de sprint y en el control del proyecto.
\end{itemize}

\Needspace{24\baselineskip}
\section{Firmas}
\noindent En señal de conformidad con el presente documento, los integrantes del equipo firman a continuación:

\vspace{1.5cm}
\begin{center}
\begin{tabular}{c c c}
    \parbox{6.5cm}{\centering \rule{6.0cm}{0.6pt}\\[5pt] \textbf{Clyde Neil Paricahua Pari}\\[2pt] {\footnotesize Project Manager}}
    & \hspace{0.8cm} &
    \parbox{6.5cm}{\centering \rule{6.0cm}{0.6pt}\\[5pt] \textbf{Ruth Karina Apaza Solis}\\[2pt] {\footnotesize Analista de Planificación}}
    \\[2.0cm]
    \parbox{6.5cm}{\centering \rule{6.0cm}{0.6pt}\\[5pt] \textbf{Eliseo Tarqui Ajahuana}\\[2pt] {\footnotesize Especialista Técnico / Arquitectura}}
    & \hspace{0.8cm} &
    \parbox{6.5cm}{\centering \rule{6.0cm}{0.6pt}\\[5pt] \textbf{Viviana Mary Quisbert Quispe}\\[2pt] {\footnotesize Gestora de Comunicación, Riesgos y Documentación}}
\end{tabular}
\end{center}
\end{document}
